\documentclass{book}

\title{Analisi Matematica T-A}
\author{Michele Mazza}
\date{November 2023}


% Pacchetti di Codifica e Lingua
\usepackage[utf8]{inputenc}
\usepackage[T1]{fontenc}
\usepackage[italian]{babel}

% Pacchetti Grafici e Diagrammi
\usepackage{graphicx}
\usepackage{tikz}
\usepackage{tikz-3dplot}
\usetikzlibrary{angles,quotes,arrows.meta, positioning}
\usepackage{pgfplots}
\pgfplotsset{compat=1.18}

% Pacchetti Matematici
\usepackage{amsmath, amssymb, amsthm}

% Pacchetti per Tabelle e Unità di Misura
\usepackage{geometry}
\usepackage{siunitx}
\usepackage{booktabs}
\usepackage{tabularx}
\usepackage{colortbl}

% Pacchetti per Layout e Formattazione
\usepackage{fancyhdr}
\usepackage{float}
\usepackage{microtype}

% Altri Pacchetti Utili
\usepackage{mdframed}
\usepackage{animate}

% Pacchetto enumitem
\usepackage{enumitem}

% Definizioni dei Colori
\definecolor{lightblue}{HTML}{9cd8fc}
\definecolor{myblue}{HTML}{139FE9}

% Definizione del comando \norm
\newcommand{\norm}[1]{\left\lVert#1\right\rVert}

%\font\nullfont=cmr10 

% Ambienti Personalizzati
\newmdtheoremenv[
    backgroundcolor=orange!10,
    linecolor=myblue,
    linewidth=0pt,
    leftmargin=0pt,
    rightmargin=0pt,
    innerleftmargin=10pt,
    innerrightmargin=10pt,
    innertopmargin=10pt,
    innerbottommargin=10pt,
    nobreak
]{curiosity}{Curiosità}

\newmdtheoremenv[
    font=\normalfont\small,
    linewidth=0pt,
    innerleftmargin=10pt,
    innerrightmargin=10pt,
]{example}{Esempio}

\newmdtheoremenv[
  backgroundcolor=lightblue,
  linecolor=myblue,
  linewidth=1.5pt,
  innertopmargin=10pt,
  innerbottommargin=10pt,
  innerrightmargin=10pt,
  innerleftmargin=10pt,
  nobreak
]{theorem}{Teorema}

\newmdtheoremenv[
  backgroundcolor=gray!10,
  linewidth=0pt,
  innertopmargin=10pt,
  innerbottommargin=10pt,
  innerrightmargin=10pt,
  innerleftmargin=10pt,
  nobreak
]{definition}{Definizione}

\newmdtheoremenv[
  backgroundcolor=gray!10,
  linecolor=myblue,
  linewidth=1.5pt,
  innertopmargin=10pt,
  innerbottommargin=10pt,
  innerrightmargin=10pt,
  innerleftmargin=10pt,
  nobreak
]{proposition}{Proposizione}

% Configurazioni di fancyhdr
\pagestyle{fancy}
\fancyhf{} % Pulisce tutti i campi
\fancyhead[LE,RO]{\thepage} % Numero di pagina
\fancyhead[RE]{\nouppercase{\leftmark}} % Titolo del capitolo
\fancyhead[LO]{\nouppercase{\rightmark}} % Titolo della sezione
\renewcommand{\headrulewidth}{0.5pt}
\renewcommand{\footrulewidth}{0pt}

\geometry{a4paper, left=20mm, right=20mm, top=20mm, bottom=20mm}

% Pacchetto Hyperref (inserito per ultimo)
\usepackage{hyperref}

\begin{document}

\frontmatter
\maketitle
\tableofcontents
\mainmatter

\chapter{Calcolo differenziale}
Nell'analisi di una funzione, il focus è sulle sue variazioni rispetto alla variabile indipendente. Il tasso di variazione, ovvero la variazione della funzione per unità di variazione della variabile, è particolarmente rivelatore. Questo concetto è fondamentale in diversi ambiti, come la fisica (per esempio, la velocità di un oggetto in movimento) o l'economia (variazioni di un indice economico).

La derivata di una funzione in un punto specifico fornisce una misura istantanea di questo tasso di variazione. Geometricamente, corrisponde al coefficiente angolare della tangente al grafico della funzione in quel punto. Il significato della derivata può variare a seconda del contesto: ad esempio, in fisica può rappresentare la velocità, mentre in economia può indicare l'utilità marginale.

Il concetto di derivata fu sviluppato in modo significativo da matematici come Leibniz e Newton, e successivamente perfezionato da Lagrange, Cauchy e Weierstrass. Il loro lavoro ha trasformato il Calcolo differenziale in uno strumento potente, utilizzato in una vasta gamma di discipline scientifiche e tecnologiche. In particolare, le equazioni differenziali, che esprimono relazioni tra le derivate di funzioni, sono fondamentali nella modellizzazione matematica di fenomeni reali.

Nel nostro studio delle derivate, affronteremo vari teoremi fondamentali del Calcolo differenziale, tra cui i teoremi di Fermat, Rolle, Lagrange, Cauchy e de l'Hôpital, esplorando le loro applicazioni. Questi teoremi aiutano a comprendere le proprietà di monotonia e i punti di estremo locale delle funzioni, nonché le relazioni tra le derivate di ordine superiore e la convessità o concavità dei grafici delle funzioni.

\section{La derivata}
\begin{definition}
    \subsection{Rapporto incrementale}
    Iniziamo introducendo il concetto di derivata di una funzione. Sia $f: \text{dom} f \subseteq \mathbb{R} \to \mathbb{R}$ una funzione reale di variabile reale; sia $x_0 \in \text{dom} f$ e supponiamo che $f$ sia definita in tutto un intorno $I_{r}(x_0)$ di $x_0$. Fissato $x \in I_{r}(x_0)$, $x \neq x_0$, indichiamo con
    \[
        \Delta x = x - x_0
    \]
    l'incremento (positivo o negativo) della variabile indipendente tra $x_0$ e $x$, e con
    \[
        \Delta f = f(x) - f(x_0)
    \]
    il corrispondente incremento della variabile dipendente. Notiamo che, dalle definizioni, segue immediatamente che $x = x_0 + \Delta x$ e $f(x) = f(x_0) + \Delta f$. Il quoziente
    \begin{equation}
        \frac{\Delta f}{\Delta x} = \frac{f(x) - f(x_0)}{x - x_0} = \frac{f(x_0 + \Delta x) - f(x_0)}{\Delta x}
    \end{equation}

    dicesi \textbf{rapporto incrementale} della funzione $f$ tra $x_0$ e $x$.
\end{definition}
\subsubsection{Incremento assoluto della variabile dipendente}
La quantità
\[\Delta f\]
rappresenta l'incremento assoluto della variabile dipendente $f$ nel passaggio da $x_0$ a $x_0 + \Delta x$.
\subsubsection{Tasso di incremento della variabile dipendente}
Il rapporto incrementale
\[
    \frac{\Delta f}{\Delta x}
\]
rappresenta il tasso di incremento.
\subsubsection{Incremento relativo della variabile dipendente}
Inoltre, la quantità
\[\frac{\Delta f}{f}\]
rappresenta l'incremento relativo di $f$.
\subsubsection{Tasso di incremento percentuale della variabile dipendente}
Se moltiplichiamo per $100$ il rapporto incrementale, otteniamo il cosiddetto tasso di incremento percentuale.

\subsubsection{Retta secante}
Dal punto di vista geometrico, il rapporto incrementale è il coefficiente angolare della retta secante al grafico di $f$ passante per i punti $P_0(x_0, f(x_0))$ e $P(x, f(x))$.
\begin{align*}
    y=s(x)=f(x_0) + \frac{\Delta f}{\Delta x}(x - x_0)
\end{align*}
\begin{figure}[H]
    \centering
    \begin{tikzpicture}[scale = 1,
            declare function={
                    f(\x) = 0.5 * \x^2 - 1.5 * \x + 2;
                    m(\p,\t) = (f(\t) - f(\p)) / (\t - \p);
                    b(\p,\t) = f(\p) - m(\p,\t) * \p;
                }]
        \def\a{0.1} %Estremi della funzione
        \def\b{4.5} %Estremi della funzione
        \def\p{1.5} %Punto 1
        \def\t{3.5} %Punto 2
        \begin{axis}[
                xlabel={\( x \)},
                ylabel={\( f(x) \)},
                xmin=-2, xmax=5,
                ymin=0, ymax=5,
                grid=both,
                axis lines=middle,
                axis line style=->,
                grid style={line width=.1pt, draw=gray!10},
                major grid style={line width=.2pt,draw=gray!50},
                minor tick num=3,
                enlargelimits={abs=0.5},
                xlabel style={at={(ticklabel* cs:1)}, anchor=north west},
                ylabel style={at={(ticklabel* cs:1)}, anchor=south west},
                xticklabels={},
                yticklabels={}
                samples=100,
                unbounded coords=jump,
                legend pos=south east
            ]
            \draw [dashed, thick, color=gray] (axis cs: \p, {f(\p)}) -- (axis cs: \p, 0);
            \draw [dashed, thick, color=gray] (axis cs: \p, {f(\p)}) -- (axis cs: 0, {f(\p)});
            \draw [dashed, thick, color=gray] (axis cs: \t, {f(\t)}) -- (axis cs: \t, 0);
            \draw [dashed, thick, color=gray] (axis cs: \t, {f(\t)}) -- (axis cs: 0, {f(\t)});
            \draw[ultra thick, color=myblue] (axis cs: \p,0) -- (axis cs: \t,0);
            \draw[ultra thick, color=myblue] (axis cs: 0,{f(\p)}) -- (axis cs: 0,{f(\t)});
            \addplot[domain=\a:\b, ultra thick, smooth, color=red]{m(\p,\t)*x + b(\p,\t)};
            \addplot[domain=\a:\b, ultra thick, smooth, color=myblue]{f(x)};
            \filldraw[color=myblue] (axis cs: \p, {f(\p)}) circle (1.5pt);
            \filldraw[color=myblue] (axis cs: \t, {f(\t)}) circle (1.5pt);
            \node[below] at (axis cs: \p, 0) {\(x_0\)};
            \node[below] at (axis cs: \t, 0) {\(x_0 + \varDelta x\)};
            \node[left, font=\small] at (axis cs: 0, {f(\p)}) {\(f(x_0)\)};
            \node[left, font=\small] at (axis cs: 0, {f(\t)}) {\(f(x_0 + \varDelta x)\)};
            \node[left, font=\small, color=red] at (axis cs: {(\p+\t)/2}, {f(\p)+f(\t)}/2) {\(s(x)\)};
            \node[above left, font=\small] at (axis cs: \p, {f(\p)}) {\(P_0\)};
            \node[right] at (axis cs: \t, {f(\t)}) {\(P\)};
        \end{axis}
    \end{tikzpicture}
    \caption{Retta secante}
\end{figure}

\noindent\begin{minipage}{\textwidth}
    \begin{definition}
        \subsection{Derivata}
        Sia $f: \text{dom} f \subseteq \mathbb{R} \to \mathbb{R}$ una funzione reale di variabile reale; sia $x_0 \in \text{dom} f$ e supponiamo che $f$ sia definita in tutto un intorno $I_{r}(x_0)$ di $x_0$. Se esiste il limite
        \begin{align}
            f'(x_0) & = \lim_{x \to x_0} \frac{f(x) - f(x_0)}{x - x_0}                    \\
                    & = \lim_{\Delta x \to 0} \frac{f(x_0 + \Delta x) - f(x_0)}{\Delta x}
        \end{align}
        esso si dice \textbf{derivata (prima)} di $f$ nel punto $x_0$ e si indica con
        \begin{align}
            f'(x_0) \quad y'(x_0) \quad \frac{df}{dx}(x_0) \quad \frac{dy}{dx}(x_0) \quad \frac{d}{dx}f(x_0) \quad \frac{d}{dx}y(x_0) \quad Df(x_0) \quad Dy(x_0)
        \end{align}
    \end{definition}
    \subsubsection{Retta tangente}
    Dal punto di vista geometrico, la derivata è il coefficiente angolare della retta tangente al grafico di $f$ nel punto $P_0(x_0, f(x_0))$.
    \begin{align*}
        y=t(x)=f(x_0) + f'(x_0)(x - x_0)
    \end{align*}
    \begin{figure}[H]
        \centering
        \begin{animateinline}[controls,loop]{5}
            \multiframe{6}{rDelta=4+-0.5}{%
                \begin{tikzpicture}[scale = 1,
                        declare function={
                                f(\x) = 0.5 * \x^2 - 1.5 * \x + 2;
                                m(\p,\t) = (f(\t) - f(\p)) / (\t - \p);
                                b(\p,\t) = f(\p) - m(\p,\t) * \p;
                            }]
                    \def\a{0.1} %Estremi della funzione
                    \def\b{4.5} %Estremi della funzione
                    \def\p{1.5} %Punto 1
                    \def\t{\rDelta} %Punto 2
                    \begin{axis}[
                            xlabel={\( x \)},
                            ylabel={\( f(x) \)},
                            xmin=-2, xmax=5,
                            ymin=0, ymax=5,
                            grid=both,
                            axis lines=middle,
                            axis line style=->,
                            grid style={line width=.1pt, draw=gray!10},
                            major grid style={line width=.2pt,draw=gray!50},
                            minor tick num=3,
                            enlargelimits={abs=0.5},
                            xlabel style={at={(ticklabel* cs:1)}, anchor=north west},
                            ylabel style={at={(ticklabel* cs:1)}, anchor=south west},
                            xticklabels={},
                            yticklabels={}
                            samples=100,
                            unbounded coords=jump,
                            legend pos=south east
                        ]
                        \draw [dashed, thick, color=gray] (axis cs: \p, {f(\p)}) -- (axis cs: \p, 0);
                        \draw [dashed, thick, color=gray] (axis cs: \p, {f(\p)}) -- (axis cs: 0, {f(\p)});
                        \draw [dashed, thick, color=gray] (axis cs: \t, {f(\t)}) -- (axis cs: \t, 0);
                        \draw [dashed, thick, color=gray] (axis cs: \t, {f(\t)}) -- (axis cs: 0, {f(\t)});
                        \draw[ultra thick, color=myblue] (axis cs: \p,0) -- (axis cs: \t,0);
                        \draw[ultra thick, color=myblue] (axis cs: 0,{f(\p)}) -- (axis cs: 0,{f(\t)});
                        \addplot[domain=\a:\b, ultra thick, smooth, color=red]{m(\p,\t)*x + b(\p,\t)};
                        \addplot[domain=\a:\b, ultra thick, smooth, color=myblue]{f(x)};
                        \filldraw[color=myblue] (axis cs: \p, {f(\p)}) circle (1.5pt);
                        \filldraw[color=myblue] (axis cs: \t, {f(\t)}) circle (1.5pt);
                        \node[left, font=\small, color=red] at (axis cs: {(\p+\t)/2}, {f(\p)+f(\t)}/2) {\(t(x)\)};
                        \node[below, font=\small, color=black] at (axis cs: {(\p+\t)/2}, 0) {\(\varDelta x \to 0\)};
                    \end{axis}
                \end{tikzpicture}
            }
        \end{animateinline}
        \begin{tikzpicture}[scale = 1,
                declare function={
                        f(\x) = 0.5 * \x^2 - 1.5 * \x + 2;
                        m(\p,\t) = (f(\t) - f(\p)) / (\t - \p);
                        b(\p,\t) = f(\p) - m(\p,\t) * \p;
                    }]
            \def\a{0.1} %Estremi della funzione
            \def\b{4.5} %Estremi della funzione
            \def\p{1.5} %Punto 1
            \def\t{1.5} %Punto 2
            \begin{axis}[
                    xlabel={\( x \)},
                    ylabel={\( f(x) \)},
                    xmin=-2, xmax=5,
                    ymin=0, ymax=5,
                    grid=both,
                    axis lines=middle,
                    axis line style=->,
                    grid style={line width=.1pt, draw=gray!10},
                    major grid style={line width=.2pt,draw=gray!50},
                    minor tick num=3,
                    enlargelimits={abs=0.5},
                    xlabel style={at={(ticklabel* cs:1)}, anchor=north west},
                    ylabel style={at={(ticklabel* cs:1)}, anchor=south west},
                    xticklabels={},
                    yticklabels={}
                    samples=100,
                    unbounded coords=jump,
                    legend pos=south east
                ]
                \draw [dashed, thick, color=gray] (axis cs: \p, {f(\p)}) -- (axis cs: \p, 0);
                \draw [dashed, thick, color=gray] (axis cs: \p, {f(\p)}) -- (axis cs: 0, {f(\p)});
                \draw [dashed, thick, color=gray] (axis cs: \t, {f(\t)}) -- (axis cs: \t, 0);
                \draw [dashed, thick, color=gray] (axis cs: \t, {f(\t)}) -- (axis cs: 0, {f(\t)});
                \draw[ultra thick, color=myblue] (axis cs: \p,0) -- (axis cs: \t,0);
                \draw[ultra thick, color=myblue] (axis cs: 0,{f(\p)}) -- (axis cs: 0,{f(\t)});
                \addplot[domain=\a:\b, ultra thick, smooth, color=red]{m(\p,\t)*x + b(\p,\t)};
                \addplot[domain=\a:\b, ultra thick, smooth, color=myblue]{f(x)};
                \filldraw[color=myblue] (axis cs: \p, {f(\p)}) circle (1.5pt);
                \filldraw[color=myblue] (axis cs: \t, {f(\t)}) circle (1.5pt);
                \node[left, font=\small, color=red] at (axis cs: {(\p+\t)/2}, {f(\p)+f(\t)}/2) {\(t(x)\)};
                \node[below, font=\small, color=black] at (axis cs: {(\p+\t)/2}, 0) {\(\varDelta x \to 0\)};
            \end{axis}
        \end{tikzpicture}
        \caption{Retta tangente}
    \end{figure}
\end{minipage}

\begin{definition}
    \subsection{Derivabilità}
    Sia $f: \text{dom} f \subseteq \mathbb{R} \to \mathbb{R}$ una funzione reale di variabile reale; sia $x_0 \in \text{dom} f$ e supponiamo che $f$ sia definita in tutto un intorno $I_{r}(x_0)$ di $x_0$. Se esiste il limite (finito) del rapporto incrementale per $\Delta x \to 0$, allora $f$ si dice \textbf{derivabile} nel punto $x_0$.
    \begin{align*}
        \lim_{\Delta x \to 0} \frac{f(x_0 + \Delta x) - f(x_0)}{\Delta x} = \lim_{x \to x_0} \frac{f(x) - f(x_0)}{x - x_0}
    \end{align*}
    Esiste il limite se e solo se esistono i limiti destro e sinistro e sono uguali.
\end{definition}

\begin{proposition}
    \subsection{Derivabilità implica continuità}
    Se $f$ è derivabile in $x_0$, allora $f$ è continua in $x_0$.
\end{proposition}

\begin{proposition}
    \subsection{Continuità non implica derivabilità}
    La funzione $f(x) = |x|$ è continua in $x_0 = 0$, ma non è derivabile in $x_0 = 0$.
\end{proposition}
\begin{figure}[H]
    \centering
    \begin{tikzpicture}[scale = 1,
            declare function={
                    f(\x) = abs(\x);
                    m(\p,\t) = (f(\t) - f(\p)) / (\t - \p);
                    b(\p,\t) = f(\p) - m(\p,\t) * \p;
                }]
        \def\a{-3} %Estremi della funzione
        \def\b{3} %Estremi della funzione
        \def\p{0} %Punto 1
        \def\t{0} %Punto 2
        \begin{axis}[
                xlabel={\( x \)},
                ylabel={\( f(x) \)},
                xmin=-3, xmax=3,
                ymin=-1, ymax=3,
                grid=both,
                axis lines=middle,
                axis line style=->,
                grid style={line width=.1pt, draw=gray!10},
                major grid style={line width=.2pt,draw=gray!50},
                minor tick num=3,
                enlargelimits={abs=0.5},
                xlabel style={at={(ticklabel* cs:1)}, anchor=north west},
                ylabel style={at={(ticklabel* cs:1)}, anchor=south west},
                xticklabels={},
                yticklabels={}
                samples=100,
                unbounded coords=jump,
                legend pos=south east
            ]
            \addplot[domain=\a:\b, ultra thick, smooth, color=myblue]{f(x)};
            \legend{\(|x|\)}
        \end{axis}
    \end{tikzpicture}
\end{figure}
La funzione è definita a tratti, in questo modo:
\begin{align*}
    f(x) = |x|\begin{cases}
                  x  & \text{se } x \geq 0 \\
                  -x & \text{se } x < 0
              \end{cases}
\end{align*}
La sua derivata in $x_0 = 0$ è:
\begin{align*}
     & \lim_{x \to 0^-} \frac{|x| - |0|}{x - 0} = \lim_{x \to 0^-} \frac{-x}{x} = -1 \\
     & \lim_{x \to 0^+} \frac{|x| - |0|}{x - 0} = \lim_{x \to 0^+} \frac{x}{x} = 1
\end{align*}
Dunque i limiti destro e sinistro non coincidono, quindi la derivata non esiste.

\section{Regole di derivazione}
\begin{theorem}
    \subsection{Algebra delle derivate}
    Di seguito l'algebra delle derivate.
    Siano $f$ e $g$ due funzioni derivabili in $x_0 \in \mathbb{R}$
    Allora sono derivabili anche le funzioni:
    \begin{center}
        \renewcommand{\arraystretch}{2} % Aumenta lo spazio tra le righe
        \begin{tabular}{l l}
            $D[f(x) \pm g(x)]                $               & $ = D[f(x)] \pm D[g(x)]                                                     $                 \\
            $D[f(x) \cdot g(x)]              $               & $ = f(x)D[g(x)] + g(x)D[f(x)]                                               $                 \\
            $D\left[\displaystyle \frac{f(x)}{g(x)}\right] $ & $ = \displaystyle \frac{g(x)D[f(x)] - f(x)D[g(x)]}{g(x)^2} \quad \text{se } g(x_0) \neq 0 $   \\
            $D[\alpha f(x) + \beta g(x)]                 $   & $ = \alpha D[f(x)] + \beta D[g(x)] \quad \text{(Linearità della derivata)}                  $ \\
            $D[f(g(x))]                      $               & $ = f'(g(x))D[g(x)] \quad \text{(Regola della catena)}$                                       \\
        \end{tabular}
    \end{center}
\end{theorem}
\begin{theorem}
    \subsection{Derivate delle funzioni elementari}
    Di seguito le derivate delle funzioni elementari.
    \begin{center}
        \renewcommand{\arraystretch}{2} % Aumenta lo spazio tra le righe
        \begin{tabular}{l l}
            $D[c] = 0 \text{ (dove $c \in \mathbb{R}$ è una costante)}$                   & $D[\alpha f(x)] = \alpha D[f(x)]$                                                                     \\
            $D[f(x) + g(x)] = D[f(x)] + D[g(x)]$                                          & $D[f(x)g(x)] = f(x)D[g(x)] + g(x)D[f(x)]$                                                             \\
            $D[f(g(x))] = f'(g(x))g'(x)$                                                  & $\displaystyle D\left[\frac{f(x)}{g(x)}\right] = \frac{g(x) D[f(x)] - f(x) D[g(x)]}{g(x)^2}$          \\
            $D[x] = 1$                                                                    & $\displaystyle D[f(x)^{g(x)}] = f(x)^{g(x)} \left[ g'(x) \ln(f(x)) + \frac{g(x) f'(x)}{f(x)} \right]$ \\
            $D[x^\alpha] = \alpha x^{\alpha-1}$                                           & $D[g(x)^\alpha] = \alpha g(x)^{\alpha-1} D[g(x)]$                                                     \\
            $D[a^x] = a^x \ln a$                                                          & $D[a^{g(x)}] = a^{g(x)} \ln a \cdot D[g(x)]$                                                          \\
            $D[e^x] = e^x$                                                                & $D[e^{g(x)}] = e^{g(x)} D[g(x)]$                                                                      \\
            $D[\ln |x|] = \displaystyle \frac{1}{x \ln (e)} = \displaystyle \frac{1}{x} $ & $D[\ln g(x)] = \displaystyle \frac{1}{g(x)} D[g(x)]$                                                  \\
            $D[\log_a |x|] = \displaystyle\frac{1}{x \ln (a)} ( \forall a>0, a\neq0 )$    & $D[\log_a g(x)] = \displaystyle \frac{1}{g(x) \ln a} D[g(x)]$                                         \\
            $D[\sin x] = \cos x$                                                          & $D[\sin g(x)] = \cos g(x) D[g(x)]$                                                                    \\
            $D[\cos x] = -\sin x$                                                         & $D[\cos g(x)] = -\sin g(x) D[g(x)]$                                                                   \\
            $D[\tan x] = 1 + \tan^2 x = \sec^2 x = \displaystyle \frac{1}{cos^2 x} $      & $D[\tan g(x)] = \sec^2 g(x) D[g(x)]$                                                                  \\
            $D[\arcsin x] = \displaystyle \frac{1}{\sqrt{1-x^2}}$                         & $D[\arcsin g(x)] = \displaystyle \frac{1}{\sqrt{1-g(x)^2}} D[g(x)]$                                   \\
            $D[\arccos x] = \displaystyle \frac{-1}{\sqrt{1-x^2}}$                        & $D[\arccos g(x)] = \displaystyle \frac{-1}{\sqrt{1-g(x)^2}} D[g(x)]$                                  \\
            $D[\arctan x] = \displaystyle \frac{1}{1+x^2}$                                & $D[\arctan g(x)] = \displaystyle \frac{1}{1+g(x)^2} D[g(x)]$                                          \\
            $D[\sinh x] = \cosh x$                                                        & $D[\sinh g(x)] = \cosh g(x) D[g(x)]$                                                                  \\
            $D[\cosh x] = \sinh x$                                                        & $D[\cosh g(x)] = \sinh g(x) D[g(x)]$                                                                  \\
            $D[\tanh x] = \text{sech}^2 x$                                                & $D[\tanh g(x)] = \text{sech}^2 g(x) D[g(x)]$                                                          \\
            $D[\text{sech}\, x] = -\text{sech}\, x \tanh x$                               & $D[\text{sech}\, g(x)] = -\text{sech}\, g(x) \tanh g(x) D[g(x)]$                                      \\
            $D[\text{csch}\, x] = -\text{csch}\, x \coth x$                               & $D[\text{csch}\, g(x)] = -\text{csch}\, g(x) \coth g(x) D[g(x)]$                                      \\
            $D[\text{coth}\, x] = -\text{csch}^2 x$                                       & $D[\text{coth}\, g(x)] = -\text{csch}^2 g(x) D[g(x)]$                                                 \\
        \end{tabular}
    \end{center}
\end{theorem}
\begin{theorem}
    \subsection{Derivate delle funzioni inverse}
    Sia \( f(x) \) una funzione continua e invertibile in un intorno di un punto \( x_0 \in \mathbb{R} \); inoltre, sia \( f \) derivabile in \( x_0 \), con \( f'(x_0) \neq 0 \). Allora la funzione inversa \( f^{-1}(y) \) è derivabile in \( y_0 = f(x_0) \) e si ha:
    \begin{equation}
        D[f^{-1}(y_0)] = \frac{1}{D[f(x_0)]} = \frac{1}{f'(x_0)} = \frac{1}{f'(f^{-1}(y_0))}
    \end{equation}
\end{theorem}
\begin{example}
    Consideriamo la funzione \( y = f(x) = \tan(x) \). La sua derivata e la funzione inversa sono date da:
    \begin{align*}
        f'(x)     & = 1 + \tan^2(x) \\
        f^{-1}(y) & = \arctan(y)
    \end{align*}

    Applicando la formula per la derivata della funzione inversa, otteniamo la derivata di \( f^{-1} \):
    \begin{align*}
        \left( f^{-1} \right)'(y) & = \frac{1}{f'(x)}         \\
                                  & = \frac{1}{1 + \tan^2(x)}
    \end{align*}

    Se ora indichiamo \( f^{-1} \) come \( g \) e torniamo a usare \( x \) come variabile indipendente, possiamo scrivere:
    \begin{align*}
        g(x)  & = \arctan(x)                       \\
        g'(x) & = \frac{1}{1 + \tan^2(\arctan(x))}
    \end{align*}

    Poiché \( \tan(\arctan(x)) = x \), la derivata finale di \( g(x) = \arctan(x) \) diventa:
    \begin{align*}
        g'(x) & = \frac{1}{1 + x^2}
    \end{align*}

    Quindi, la derivata di \( q(x) = \arctan(x) \) è:
    \begin{align*}
        q'(x) & = \frac{1}{1 + x^2}
    \end{align*}
\end{example}

\begin{definition}
    \subsection{Derivata logaritmica}
    Se abbiamo una funzione \( f(x) \) che è sempre positiva e che cambia con \( x \), possiamo guardare alla sua "derivata logaritmica". Questo è un modo di misurare quanto velocemente cambia \( f(x) \), ma in termini percentuali. Se definiamo \( g(x) = \log f(x) \), allora la sua derivata, che è la derivata logaritmica di \( f \), si calcola così:
    \begin{align*}
        g'(x) = \frac{f'(x)}{f(x)}
    \end{align*}

    In pratica, questa derivata logaritmica ci dice quanto rapidamente \( f(x) \) sta crescendo o diminuendo, rispetto al suo valore attuale. È particolarmente utile quando \( f(x) \) varia molto, perché ci dà una visione chiara di quanto grandi siano queste variazioni in proporzione.
\end{definition}

\section{Punti di non derivabilità}
\begin{definition}
    \subsection{Derivabilità}
    Una funzione \( f \) definita in un intorno di un punto \( x_0 \in \mathbb{R} \)
    è derivabile in \( x_0 \) se e solo se è derivabile da destra e da sinistra in \( x_0 \) e le
    derivate destra e sinistra coincidono. In tal caso si ha
    \begin{align*}
        f'(x_0) & \quad = \quad f'_+(x_0) \quad = \quad f'_-(x_0)                                                                               \\
                & \quad=  \quad \lim_{x \to x_0^+} \frac{f(x) - f(x_0)}{x - x_0} \quad = \quad \lim_{x \to x_0^-} \frac{f(x) - f(x_0)}{x - x_0}
    \end{align*}
    Tuttavia, per garantire la derivabilità di \( f \) in \( x_0 \), è necessario verificare anche l'ipotesi di continuità nel punto \( x_0 \).
\end{definition}

\noindent\begin{minipage}{\textwidth}
    \begin{definition}
        \subsection{Punto angoloso}
        Un punto \( x_0 \) è detto \textbf{angoloso} se la derivata destra e sinistra esistono, ma non coincidono.
        \begin{align*}
            f'_+(x_0) \neq f'_-(x_0)
        \end{align*}
    \end{definition}
    \begin{figure}[H]
        \centering
        \begin{animateinline}[controls,loop]{1}
            \multiframe{2}{rDelta=-0.1+0.2}{%
                \begin{tikzpicture}[scale = 1,
                        declare function={
                                f(\x) = abs(\x);
                                m(\p,\t) = (f(\t) - f(\p)) / (\t - \p);
                                b(\p,\t) = f(\p) - m(\p,\t) * \p;
                            }]
                    \def\a{-4.5} %Estremi della funzione
                    \def\b{4.5} %Estremi della funzione
                    \def\p{0} %Punto 1
                    \def\t{\rDelta} %Punto 2
                    \begin{axis}[
                            title={Punto \textbf{angoloso}},
                            title style={at={(0.5,1.2)}, anchor=north},
                            xlabel={\( x \)},
                            ylabel={\( f(x) \)},
                            xmin=-5, xmax=5,
                            ymin=-5, ymax=5,
                            restrict y to domain=-10:10,
                            grid=both,
                            axis lines=middle,
                            axis line style=->,
                            grid style={line width=.1pt, draw=gray!10},
                            major grid style={line width=.2pt,draw=gray!50},
                            minor tick num=3,
                            enlargelimits={abs=0.5},
                            xlabel style={at={(ticklabel* cs:1)}, anchor=north west},
                            ylabel style={at={(ticklabel* cs:1)}, anchor=south west},
                            xticklabels={},
                            yticklabels={}
                            samples=100,
                            unbounded coords=jump,
                            legend pos=south east
                        ]
                        \draw [dashed, thick, color=gray] (axis cs: \p, {f(\p)}) -- (axis cs: \p, 0);
                        \draw [dashed, thick, color=gray] (axis cs: \p, {f(\p)}) -- (axis cs: 0, {f(\p)});
                        \draw [dashed, thick, color=gray] (axis cs: \t, {f(\t)}) -- (axis cs: \t, 0);
                        \draw [dashed, thick, color=gray] (axis cs: \t, {f(\t)}) -- (axis cs: 0, {f(\t)});
                        \draw[ultra thick, color=myblue] (axis cs: \p,0) -- (axis cs: \t,0);
                        \draw[ultra thick, color=myblue] (axis cs: 0,{f(\p)}) -- (axis cs: 0,{f(\t)});
                        \addplot[domain=\a:\b, ultra thick, smooth, color=red]{m(\p,\t)*x + b(\p,\t)};
                        \addplot[domain=\a:\b, ultra thick, smooth, color=myblue]{f(x)};
                        \filldraw[color=myblue] (axis cs: \p, {f(\p)}) circle (1.5pt);
                        \filldraw[color=myblue] (axis cs: \t, {f(\t)}) circle (1.5pt);
                        \node[left, font=\small, color=red] at (axis cs: {(\p+\t)/2}, {f(\p)+f(\t)}/2) {\(t(x)\)};
                        \node[below, font=\small, color=black] at (axis cs: {(\p+\t)/2}, 0) {\(\varDelta x \to 0\)};
                        \legend{\(t(x)\),\(|x|\)}
                    \end{axis}
                \end{tikzpicture}
            }
        \end{animateinline}
        \begin{tikzpicture}[scale = 1,
                declare function={
                        f(\x) = abs(\x);
                        m(\p,\t) = (f(\t) - f(\p)) / (\t - \p);
                        b(\p,\t) = f(\p) - m(\p,\t) * \p;
                    }]
            \def\a{-4.5} %Estremi della funzione
            \def\b{4.5} %Estremi della funzione
            \def\p{0} %Punto 1
            \def\t{+0.1} %Punto 2
            \begin{axis}[
                    title={Punto \textbf{angoloso}},
                    title style={at={(0.5,1.2)}, anchor=north},
                    xlabel={\( x \)},
                    ylabel={\( f(x) \)},
                    xmin=-5, xmax=5,
                    ymin=-5, ymax=5,
                    restrict y to domain=-10:10,
                    grid=both,
                    axis lines=middle,
                    axis line style=->,
                    grid style={line width=.1pt, draw=gray!10},
                    major grid style={line width=.2pt,draw=gray!50},
                    minor tick num=3,
                    enlargelimits={abs=0.5},
                    xlabel style={at={(ticklabel* cs:1)}, anchor=north west},
                    ylabel style={at={(ticklabel* cs:1)}, anchor=south west},
                    xticklabels={},
                    yticklabels={}
                    samples=100,
                    unbounded coords=jump,
                    legend pos=south east
                ]
                \draw [dashed, thick, color=gray] (axis cs: \p, {f(\p)}) -- (axis cs: \p, 0);
                \draw [dashed, thick, color=gray] (axis cs: \p, {f(\p)}) -- (axis cs: 0, {f(\p)});
                \draw [dashed, thick, color=gray] (axis cs: \t, {f(\t)}) -- (axis cs: \t, 0);
                \draw [dashed, thick, color=gray] (axis cs: \t, {f(\t)}) -- (axis cs: 0, {f(\t)});
                \draw[ultra thick, color=myblue] (axis cs: \p,0) -- (axis cs: \t,0);
                \draw[ultra thick, color=myblue] (axis cs: 0,{f(\p)}) -- (axis cs: 0,{f(\t)});
                \addplot[domain=\a:\b, ultra thick, smooth, color=red]{m(\p,\t)*x + b(\p,\t)};
                \addplot[domain=\a:\b, ultra thick, smooth, color=myblue]{f(x)};
                \filldraw[color=myblue] (axis cs: \p, {f(\p)}) circle (1.5pt);
                \filldraw[color=myblue] (axis cs: \t, {f(\t)}) circle (1.5pt);
                \node[left, font=\small, color=red] at (axis cs: {(\p+\t)/2}, {f(\p)+f(\t)}/2) {\(t(x)\)};
                \node[below, font=\small, color=black] at (axis cs: {(\p+\t)/2}, 0) {\(\varDelta x \to 0\)};
                \legend{\(t(x)\),\(|x|\)}
            \end{axis}
        \end{tikzpicture}
        \caption{Punto angoloso \( f'_+(x_0) \neq f'_-(x_0)\)}
    \end{figure}
\end{minipage}

\noindent\begin{minipage}{\textwidth}
    \begin{definition}
        \subsection{Punto di cuspide}
        Un punto \( x_0 \) è detto \textbf{di cuspide} se la derivata destra e sinistra esistono, ma non coincidono e sono infinite.
        \begin{align*}
            f'_+(x_0) = \pm \infty \quad \text{e} \quad f'_-(x_0) = \mp \infty
        \end{align*}
    \end{definition}
    \begin{figure}[H]
        \centering
        \begin{animateinline}[controls,loop]{1}
            \multiframe{2}{rDelta=-0.005+0.010}{%
                \begin{tikzpicture}[scale = 1,
                        declare function={
                                f(\x) = abs(\x)^(1/2);
                                m(\p,\t) = (f(\t) - f(\p)) / (\t - \p);
                                b(\p,\t) = f(\p) - m(\p,\t) * \p;
                            }]
                    \def\a{-2} %Estremi della funzione
                    \def\b{2} %Estremi della funzione
                    \def\p{0} %Punto 1
                    \def\t{\rDelta} %Punto 2
                    \begin{axis}[
                            title={Punto di \textbf{cuspide}},
                            title style={at={(0.5,1.2)}, anchor=north},
                            xlabel={\( x \)},
                            ylabel={\( f(x) \)},
                            xmin=-2, xmax=2,
                            ymin=-2, ymax=2,
                            restrict y to domain=-10:10,
                            grid=both,
                            axis lines=middle,
                            axis line style=->,
                            grid style={line width=.1pt, draw=gray!10},
                            major grid style={line width=.2pt,draw=gray!50},
                            minor tick num=3,
                            enlargelimits={abs=0.5},
                            xlabel style={at={(ticklabel* cs:1)}, anchor=north west},
                            ylabel style={at={(ticklabel* cs:1)}, anchor=south west},
                            xticklabels={},
                            yticklabels={}
                            samples=100,
                            unbounded coords=jump,
                            legend pos=south east
                        ]
                        \draw [dashed, thick, color=gray] (axis cs: \p, {f(\p)}) -- (axis cs: \p, 0);
                        \draw [dashed, thick, color=gray] (axis cs: \p, {f(\p)}) -- (axis cs: 0, {f(\p)});
                        \draw [dashed, thick, color=gray] (axis cs: \t, {f(\t)}) -- (axis cs: \t, 0);
                        \draw [dashed, thick, color=gray] (axis cs: \t, {f(\t)}) -- (axis cs: 0, {f(\t)});
                        \addplot[domain=\a:\b, ultra thick, smooth, color=red]{m(\p,\t)*x + b(\p,\t)};
                        \addplot[domain=\a:\b, ultra thick, smooth, color=myblue]{f(x)};
                        \filldraw[color=myblue] (axis cs: \p, {f(\p)}) circle (1.5pt);
                        \filldraw[color=myblue] (axis cs: \t, {f(\t)}) circle (1.5pt);
                        \node[left, font=\small, color=red] at (axis cs: {(\p+\t)/2}, {f(\p)+f(\t)}/2) {\(t(x)\)};
                        \node[below right, font=\small, color=black] at (axis cs: {(\p+\t)/2}, 0) {\(\varDelta x \to 0\)};
                        \legend{\(t(x)\),\(\sqrt{|x|}\)}
                    \end{axis}
                \end{tikzpicture}
            }
        \end{animateinline}
        \begin{tikzpicture}[scale = 1,
                declare function={
                        f(\x) = abs(\x)^(1/2);
                        m(\p,\t) = (f(\t) - f(\p)) / (\t - \p);
                        b(\p,\t) = f(\p) - m(\p,\t) * \p;
                    }]
            \def\a{-2} %Estremi della funzione
            \def\b{2} %Estremi della funzione
            \def\p{0} %Punto 1
            \def\t{0.005} %Punto 2
            \begin{axis}[
                    title={Punto di \textbf{cuspide}},
                    title style={at={(0.5,1.2)}, anchor=north},
                    xlabel={\( x \)},
                    ylabel={\( f(x) \)},
                    xmin=-2, xmax=2,
                    ymin=-2, ymax=2,
                    restrict y to domain=-10:10,
                    grid=both,
                    axis lines=middle,
                    axis line style=->,
                    grid style={line width=.1pt, draw=gray!10},
                    major grid style={line width=.2pt,draw=gray!50},
                    minor tick num=3,
                    enlargelimits={abs=0.5},
                    xlabel style={at={(ticklabel* cs:1)}, anchor=north west},
                    ylabel style={at={(ticklabel* cs:1)}, anchor=south west},
                    xticklabels={},
                    yticklabels={}
                    samples=100,
                    unbounded coords=jump,
                    legend pos=south east
                ]
                \draw [dashed, thick, color=gray] (axis cs: \p, {f(\p)}) -- (axis cs: \p, 0);
                \draw [dashed, thick, color=gray] (axis cs: \p, {f(\p)}) -- (axis cs: 0, {f(\p)});
                \draw [dashed, thick, color=gray] (axis cs: \t, {f(\t)}) -- (axis cs: \t, 0);
                \draw [dashed, thick, color=gray] (axis cs: \t, {f(\t)}) -- (axis cs: 0, {f(\t)});
                \addplot[domain=\a:\b, ultra thick, smooth, color=red]{m(\p,\t)*x + b(\p,\t)};
                \addplot[domain=\a:\b, ultra thick, smooth, color=myblue]{f(x)};
                \filldraw[color=myblue] (axis cs: \p, {f(\p)}) circle (1.5pt);
                \filldraw[color=myblue] (axis cs: \t, {f(\t)}) circle (1.5pt);
                \node[left, font=\small, color=red] at (axis cs: {(\p+\t)/2}, {f(\p)+f(\t)}/2) {\(t(x)\)};
                \node[below right, font=\small, color=black] at (axis cs: {(\p+\t)/2}, 0) {\(\varDelta x \to 0\)};
                \legend{\(t(x)\),\(\sqrt{|x|}\)}
            \end{axis}
        \end{tikzpicture}
        \caption{Punto di cuspide \( f'_+(x_0) = \pm \infty \quad \text{e} \quad f'_-(x_0) = \mp \infty\)}
    \end{figure}
\end{minipage}

\noindent\begin{minipage}{\textwidth}
    \begin{definition}
        \subsection{Punto di flesso a tangente verticale}
        Un punto \( x_0 \) è detto \textbf{di flesso a tangente verticale} se la derivata destra e sinistra esistono, coincidono e sono infinite dello stesso segno.
        \begin{align*}
            f'_+(x_0) = \pm \infty \quad \text{e} \quad f'_-(x_0) = \pm \infty
        \end{align*}
    \end{definition}
    \begin{figure}[H]
        \centering
        \begin{animateinline}[controls,loop]{1}
            \multiframe{2}{rDelta=-0.005+0.010}{%
                \begin{tikzpicture}[scale = 1,
                        declare function={
                                f(\x) = (\x<0 ? -1 : 1)*abs(\x)^(1/3);
                                m(\p,\t) = (f(\t) - f(\p)) / (\t - \p);
                                b(\p,\t) = f(\p) - m(\p,\t) * \p;
                            }]
                    \def\a{-2} %Estremi della funzione
                    \def\b{2} %Estremi della funzione
                    \def\p{0} %Punto 1
                    \def\t{\rDelta} %Punto 2
                    \begin{axis}[
                            title={Punto \textbf{flesso a tangente verticale}},
                            title style={at={(0.5,1.2)}, anchor=north},
                            xlabel={\( x \)},
                            ylabel={\( f(x) \)},
                            xmin=-2, xmax=2,
                            ymin=-2, ymax=2,
                            restrict y to domain=-10:10,
                            grid=both,
                            axis lines=middle,
                            axis line style=->,
                            grid style={line width=.1pt, draw=gray!10},
                            major grid style={line width=.2pt,draw=gray!50},
                            minor tick num=3,
                            enlargelimits={abs=0.5},
                            xlabel style={at={(ticklabel* cs:1)}, anchor=north west},
                            ylabel style={at={(ticklabel* cs:1)}, anchor=south west},
                            xticklabels={},
                            yticklabels={}
                            samples=100,
                            unbounded coords=jump,
                            legend pos=south east
                        ]
                        \draw [dashed, thick, color=gray] (axis cs: \p, {f(\p)}) -- (axis cs: \p, 0);
                        \draw [dashed, thick, color=gray] (axis cs: \p, {f(\p)}) -- (axis cs: 0, {f(\p)});
                        \draw [dashed, thick, color=gray] (axis cs: \t, {f(\t)}) -- (axis cs: \t, 0);
                        \draw [dashed, thick, color=gray] (axis cs: \t, {f(\t)}) -- (axis cs: 0, {f(\t)});
                        \addplot[domain=\a:\b, ultra thick, smooth, color=red]{m(\p,\t)*x + b(\p,\t)};
                        \addplot[domain=\a:\b, ultra thick, smooth, color=myblue]{f(x)};
                        \filldraw[color=myblue] (axis cs: \p, {f(\p)}) circle (1.5pt);
                        \filldraw[color=myblue] (axis cs: \t, {f(\t)}) circle (1.5pt);
                        \node[left, font=\small, color=red] at (axis cs: {(\p+\t)/2}, {f(\p)+f(\t)}/2) {\(t(x)\)};
                        \node[below right, font=\small, color=black] at (axis cs: {(\p+\t)/2}, 0) {\(\varDelta x \to 0\)};
                        \legend{\(t(x)\),\(\sqrt[3]{x}\)}
                    \end{axis}
                \end{tikzpicture}
            }
        \end{animateinline}
        \begin{tikzpicture}[scale = 1,
                declare function={
                        f(\x) = (\x<0 ? -1 : 1)*abs(\x)^(1/3);
                        m(\p,\t) = (f(\t) - f(\p)) / (\t - \p);
                        b(\p,\t) = f(\p) - m(\p,\t) * \p;
                    }]
            \def\a{-2} %Estremi della funzione
            \def\b{2} %Estremi della funzione
            \def\p{0} %Punto 1
            \def\t{0.005} %Punto 2
            \begin{axis}[
                    title={Punto \textbf{flesso a tangente verticale}},
                    title style={at={(0.5,1.2)}, anchor=north},
                    xlabel={\( x \)},
                    ylabel={\( f(x) \)},
                    xmin=-2, xmax=2,
                    ymin=-2, ymax=2,
                    restrict y to domain=-10:10,
                    grid=both,
                    axis lines=middle,
                    axis line style=->,
                    grid style={line width=.1pt, draw=gray!10},
                    major grid style={line width=.2pt,draw=gray!50},
                    minor tick num=3,
                    enlargelimits={abs=0.5},
                    xlabel style={at={(ticklabel* cs:1)}, anchor=north west},
                    ylabel style={at={(ticklabel* cs:1)}, anchor=south west},
                    xticklabels={},
                    yticklabels={}
                    samples=100,
                    unbounded coords=jump,
                    legend pos=south east
                ]
                \draw [dashed, thick, color=gray] (axis cs: \p, {f(\p)}) -- (axis cs: \p, 0);
                \draw [dashed, thick, color=gray] (axis cs: \p, {f(\p)}) -- (axis cs: 0, {f(\p)});
                \draw [dashed, thick, color=gray] (axis cs: \t, {f(\t)}) -- (axis cs: \t, 0);
                \draw [dashed, thick, color=gray] (axis cs: \t, {f(\t)}) -- (axis cs: 0, {f(\t)});
                \addplot[domain=\a:\b, ultra thick, smooth, color=red]{m(\p,\t)*x + b(\p,\t)};
                \addplot[domain=\a:\b, ultra thick, smooth, color=myblue]{f(x)};
                \filldraw[color=myblue] (axis cs: \p, {f(\p)}) circle (1.5pt);
                \filldraw[color=myblue] (axis cs: \t, {f(\t)}) circle (1.5pt);
                \node[left, font=\small, color=red] at (axis cs: {(\p+\t)/2}, {f(\p)+f(\t)}/2) {\(t(x)\)};
                \node[below right, font=\small, color=black] at (axis cs: {(\p+\t)/2}, 0) {\(\varDelta x \to 0\)};
                \legend{\(t(x)\), \(\sqrt[3]{x}\)}
            \end{axis}
        \end{tikzpicture}
        \caption{Punto di flesso a tangente verticale \( f'_+(x_0) = \pm \infty \quad \text{e} \quad f'_-(x_0) = \pm \infty \)}
    \end{figure}
\end{minipage}

\noindent\begin{minipage}{\textwidth}
    \section{Punti di estremo e punti critici di una funzione}

    \begin{definition}
        \subsection{Punto di Massimo Relativo}
        Un punto \( x_0 \) è detto \textbf{di massimo relativo} per una funzione \( f \) se esiste un intorno \( I_r(x_0) \) tale che, per ogni \( x \in I_r(x_0) \cap \text{dom}(f) \), si ha
        \begin{align*}
            f(x) \leq f(x_0).
        \end{align*}
        Il valore \( f(x_0) \) si dice massimo relativo di \( f \).
    \end{definition}

    \begin{definition}
        \subsection{Punto di Massimo Assoluto}
        Un punto \( x_0 \) è detto \textbf{di massimo assoluto} per una funzione \( f \) se, per ogni \( x \in \text{dom}(f) \), si ha
        \begin{align*}
            f(x) \leq f(x_0).
        \end{align*}
        Il valore \( f(x_0) \) si dice massimo assoluto di \( f \). Il massimo si definisce \textbf{stretto} se
        \begin{align*}
            f(x) < f(x_0) \quad \text{per ogni } x \neq x_0.
        \end{align*}
    \end{definition}

    \begin{definition}
        \subsection{Punto di critico}
        Un punto \( x_0 \) si dice \textbf{punto critico} (o \textbf{stazionario}) di una funzione \( f \) se \( f \) è derivabile in \( x_0 \) e si ha \( f'(x_0) = 0 \).
    \end{definition}

    \begin{figure}[H]
        \centering
        \begin{tikzpicture}[scale = 1,
            declare function={
            f(\x) = {\x^3 - 2*\x^2 - \x + 2};
            m(\p,\t) = (f(\t) - f(\p)) / (\t - \p);
            b(\p,\t) = f(\p) - m(\p,\t) * \p;
            }]
            \def\a{-5} %Estremi della funzione
            \def\b{5} %Estremi della funzione
            \def\p{-0.22} %Punto 1
            \def\t{-0.22} %Punto 2
            \def\r{1.55} %Punto 3
            \begin{axis}[
                    title={Punti \textbf{critici}},
                    title style={at={(0.5,1.2)}, anchor=north},
                    xlabel={\( x \)},
                    ylabel={\( f(x) \)},
                    xmin=-5, xmax=5,
                    ymin=-4, ymax=4,
                    restrict y to domain=-10:10,
                    grid=both,
                    axis lines=middle,
                    axis line style=->,
                    grid style={line width=.1pt, draw=gray!10},
                    major grid style={line width=.2pt,draw=gray!50},
                    minor tick num=3,
                    enlargelimits={abs=0.5},
                    xlabel style={at={(ticklabel* cs:1)}, anchor=north west},
                    ylabel style={at={(ticklabel* cs:1)}, anchor=south west},
                    xticklabels={},
                    yticklabels={}
                    samples=100,
                    unbounded coords=jump,
                    legend pos=south east
                ]
                \draw [dashed, thick, color=gray] (axis cs: \p, {f(\p)}) -- (axis cs: \p, 0);
                \draw [dashed, thick, color=gray] (axis cs: \p, {f(\p)}) -- (axis cs: 0, {f(\p)});
                \draw [dashed, thick, color=gray] (axis cs: \r, {f(\r)}) -- (axis cs: \r, 0);
                \draw [dashed, thick, color=gray] (axis cs: \r, {f(\r)}) -- (axis cs: 0, {f(\r)});
                \addplot[domain=-1.5:1, ultra thick, smooth, color=red]{m(\p,\t)*x + b(\p,\t)};
                \addplot[domain=0.1:3, ultra thick, smooth, color=red]{m(\r,\r)*x + b(\r,\r)};
                \addplot[domain=\a:\b, ultra thick, smooth, color=myblue]{f(x)};
                \filldraw[color=myblue] (axis cs: \p, {f(\p)}) circle (1.5pt);
                \filldraw[color=myblue] (axis cs: \r, {f(\r)}) circle (1.5pt);
                \node[font=\small, color=myblue] at (axis cs: {(\p+\p)/2}, {f(\p)+f(\p)}/2) {\(P_1\)};
                \node[font=\small, color=myblue] at (axis cs: {(\r+\r)/2}, {f(\r)+f(\r)}/2) {\(P_2\)};
            \end{axis}
        \end{tikzpicture}
        \caption{Punti critici \( f'(x_0) = 0 \)}
    \end{figure}


    \begin{theorem}
        \subsection{Teorema di Fermat}
        Sia \( f \) una funzione derivabile in un intorno di un punto \( x_0 \). Se \( x_0 \) è un punto di massimo o di minimo per \( f \), allora
        \[f'(x_0)=0\]
        cioè, \( x_0 \) è un punto critico per \( f \).
    \end{theorem}
\end{minipage}

\noindent\begin{minipage}{\textwidth}
    \section{I Teoremi di Rolle, Lagrange e Cauchy}
    \begin{theorem}
        \subsection{Teorema di Rolle}
        Sia \(f\) definita su \([a,b]\), se:
        \begin{align*}
             & f \text{ è continua in } [a,b]   \\
             & f \text{ è derivabile in } (a,b) \\
             & f(a)=f(b)
        \end{align*}
        allora esiste \textbf{almeno un punto} \(c \in (a,b)\) tale che \(f'(c)=0\).
    \end{theorem}
    \begin{figure}[H]
        \centering
        \begin{tikzpicture}[scale = 1,
            declare function={
            f(\x) = {\x^3 - 2*\x^2 - \x + 4};
            m(\p,\t) = (f(\t) - f(\p)) / (\t - \p);
            b(\p,\t) = f(\p) - m(\p,\t) * \p;
            }]
            \def\a{2} %Estremi della funzione
            \def\b{5} %Estremi della funzione
            \def\p{2.78} %Punto 1
            \def\t{2.78} %Punto 2
            \def\r{4.55} %Punto 3
            \begin{axis}[
                    title={Teorema di \textbf{Rolle}},
                    title style={at={(0.5,1.2)}, anchor=north},
                    xlabel={\( x \)},
                    ylabel={\( f(x) \)},
                    xmin=0, xmax=5,
                    ymin=-1, ymax=4.5,
                    restrict y to domain=-10:10,
                    grid=both,
                    axis lines=middle,
                    axis line style=->,
                    grid style={line width=.1pt, draw=gray!10},
                    major grid style={line width=.2pt,draw=gray!50},
                    minor tick num=3,
                    enlargelimits={abs=0.5},
                    xlabel style={at={(ticklabel* cs:1)}, anchor=north west},
                    ylabel style={at={(ticklabel* cs:1)}, anchor=south west},
                    samples=100,
                    unbounded coords=jump,
                    legend pos=south east
                ]
                \addplot[domain=\a:\b, ultra thick, smooth, color=myblue]{f(x-3)};
                \filldraw[color=myblue] (axis cs: \p, {f(\p-3)}) circle (1.5pt);
                \filldraw[color=myblue] (axis cs: \r, {f(\r-3)}) circle (1.5pt);
                \filldraw[color=myblue] (axis cs: \a, {f(\a-3)}) circle (1.5pt);
                \filldraw[color=myblue] (axis cs: \b, {f(\b-3)}) circle (1.5pt);
                \node[above, font=\small, color=myblue] at (axis cs: {(\p)}, {f(\p)}) {\(M\)};
                \node[below, font=\small, color=myblue] at (axis cs: {(\r)}, {f(\r)}) {\(m\)};
                \node [left, font=\small] at (axis cs: {\a}, {f(\a)}) {\(f(a)\)};
                \node [right, font=\small] at (axis cs: {\b}, {f(\b-3)}) {\(f(b)\)};
                \node [above, font=\small] at (axis cs: {\p}, {f(\p-3)}) {\(M\)};
                \node [below, font=\small] at (axis cs: {\r}, {f(\r-3)}) {\(m\)};
            \end{axis}
        \end{tikzpicture}
        \caption{Teorema di Rolle}
    \end{figure}
\end{minipage}

\begin{theorem}
    \subsection{Teorema di Lagrange}
    Sia \(f\) definita su \([a,b]\), se:
    \begin{align*}
         & f \text{ è continua in } [a,b]   \\
         & f \text{ è derivabile in } (a,b)
    \end{align*}
    allora esiste \textbf{almeno un punto} \(c \in (a,b)\) tale che
    \begin{align*}
        f'(c)=\frac{f(b)-f(a)}{b-a}
    \end{align*}
    ovvero la retta tangente in \(c\) è parallela alla retta passante per i punti \((a,f(a))\) e \((b,f(b))\).
\end{theorem}

\begin{figure}[H]
    \centering
    \begin{tikzpicture}[scale = 1,
            declare function={
                    f(\x) =  2 + \x^3/3 - \x^2/2 - \x;
                    m(\p,\t) = (f(\t) - f(\p)) / (\t - \p);
                    b(\p,\t) = f(\p) - m(\p,\t) * \p;
                }]
        \def\a{-1.86} %Estremi della funzione
        \def\b{3} %Estremi della funzione
        \def\p{-1.86} %Punto 1
        \def\t{3} %Punto 2
        \begin{axis}[
                title={Teorema di \textbf{Lagrange}},
                title style={at={(0.5,1.2)}, anchor=north},
                xlabel={\( x \)},
                ylabel={\( f(x) \)},
                xmin=-4, xmax=4,
                ymin=-1, ymax=4,
                grid=both,
                axis lines=middle,
                axis line style=->,
                grid style={line width=.1pt, draw=gray!10},
                major grid style={line width=.2pt,draw=gray!50},
                minor tick num=3,
                enlargelimits={abs=0.5},
                xlabel style={at={(ticklabel* cs:1)}, anchor=north west},
                ylabel style={at={(ticklabel* cs:1)}, anchor=south west},
                xticklabels={},
                yticklabels={}
                samples=100,
                unbounded coords=jump,
                legend pos=south east
            ]
            \addplot[domain=\a:\b, thick, smooth, color=red]{m(\p,\t)*x + b(\p,\t)};
            \addplot[domain=1:\b, thick, smooth, color=red]{m(\p,\t)*x - 0.83};
            \addplot[domain=\a:0, thick, smooth, color=red]{m(\p,\t)*x +2.96};
            \addplot[domain=\a:\b, ultra thick, smooth, color=myblue]{f(x)};
            \filldraw[color=myblue] (axis cs: \p, {f(\p)}) circle (1.5pt);
            \filldraw[color=myblue] (axis cs: \t, {f(\t)}) circle (1.5pt);
            \node[left, font=\small, color=red] at (axis cs: {(\p+\t)/2}, {f(\p)+f(\t)}/2) {\(s(x)\)};
        \end{axis}
    \end{tikzpicture}
    \caption{Teorema di Lagrange}
\end{figure}

\begin{theorem}
    \subsection{Teorema di Cauchy}
    Siano \(f\) e \(g\) definite su \([a,b]\), se:
    \begin{align*}
         & f \text{ e } g \text{ sono continue in } [a,b]   \\
         & f \text{ e } g \text{ sono derivabili in } (a,b) \\
         & g'(x) \neq 0 \text{ per ogni } x \in (a,b)
    \end{align*}
    allora esiste \textbf{almeno un punto} \(c \in (a,b)\) tale che
    \begin{align*}
        \frac{f'(c)}{g'(c)}=\frac{f(b)-f(a)}{g(b)-g(a)}
    \end{align*}
    ovvero la retta tangente in \(c\) è parallela alla retta passante per i punti \((a,f(a))\) e \((b,f(b))\).
    \'E dunque un estensione del teorema di Lagrange.
\end{theorem}

\section{Intervalli di monotonia di una funzione}
I seguenti teoremi e proposizioni permettono di ricondurre la ricerca dei
punti di estremo di \( f \) e dei suoi intervalli di monotonia alla ricerca degli zeri di
\( f' \) e allo studio del suo segno.

\begin{proposition}
    Sia \(I\) un intervallo e \(f\) una funzione derivabile in \(I\). Allora:
    \begin{align*}
         & \text{se } f'(x) \geq 0 \text{ per ogni } x \in I \text{ allora } f \text{ è crescente in } I   \\
         & \text{se } f'(x) \leq 0 \text{ per ogni } x \in I \text{ allora } f \text{ è decrescente in } I
    \end{align*}
\end{proposition}

\begin{theorem}
    \subsection{Massimi e Minimi}
    Sia \( f \) derivabile sull'intervallo \( I \). Sia \( x_0 \) un punto critico
    di \( f \) interno a \( I \). Se \( f'(x) > 0 \) a sinistra di \( x_0 \) e \( f'(x) < 0 \) a destra di \( x_0 \), allora
    \( x_0 \) è punto di massimo per \( f \); viceversa, se \( f'(x) < 0 \) a sinistra di \( x_0 \) e
    \( f'(x) > 0 \) a destra di \( x_0 \), allora \( x_0 \) è punto di minimo per \( f \).
\end{theorem}

\section{Derivate di ordine superiore}
Sia \( f \) una funzione derivabile \( n \) volte in un intorno di \( x_0 \). Si definisce
\begin{definition}
    \subsection{Derivata n-esima}
    \begin{align*}
        f^{(n)}(x_0) = \left(f^{(n-1)}\right)'(x_0)
    \end{align*}
    altre notazioni sono
    \begin{align*}
        y^{(n)}(x_0) = f^{(n)}(x_0) = \frac{\mathrm{d}^n f}{\mathrm{d}x^n}(x_0) = \frac{\mathrm{d}^n}{\mathrm{d}x^n}f(x_0) = D^n f(x_0)
    \end{align*}
\end{definition}

\begin{definition}
    \subsection{Classe di derivabilità}
    Si dice che \( f \) è di \textbf{classe} \( C^n \) in un intervallo \( I \) se \( f \) è derivabile \( n \) volte in \( I \) e se \( f^{(n)} \) è continua in \( I \).
    Una funzione è di \textbf{classe} \( C^{\infty} \) se è di classe \( C^n \) per ogni \( n \in \mathbb{N} \), ad esempio tutte le funzioni elementari sul loro dominio sono di classe \( C^{\infty} \).
\end{definition}
\noindent\begin{minipage}{\textwidth}
    \section{Studio di funzione}
    \begin{proposition}
        Quando si studia una funzione \(f\) si devono seguire i seguenti passi:
        \begin{align*}
             & \text{1. } \text{Dominio di \(f\)}          \\
             & \text{2. } \text{Simmetrie di \(f\)}        \\
             & \text{3. } \text{Intersezioni con gli assi} \\
             & \text{4. } \text{Asintoti}                  \\
             & \text{5. } \text{Segno di \(f\)}            \\
             & \text{6. } \text{Monotonia di \(f\)}        \\
             & \text{7. } \text{Punti di massimo e minimo} \\
             & \text{8. } \text{Concavità e convessità}    \\
             & \text{9. } \text{Punti di flesso}           \\
             & \text{10. } \text{Grafico di \(f\)}
        \end{align*}
    \end{proposition}
\end{minipage}

\section{Teorema di de l'Hôpital}
\begin{theorem}
    Siano \(f\) e \(g\) due funzioni definite in un intorno destro di \(a\), escluso \(a\) stesso, e sia \(g'(x) \neq 0\) in tale intorno. Se
    \begin{align*}
        \lim_{x \to a} f(x) = \lim_{x \to a} g(x) = 0
    \end{align*}
    o
    \begin{align*}
        \lim_{x \to a} f(x) = \lim_{x \to a} g(x) = \pm \infty
    \end{align*}
    e se esiste il limite
    \begin{align*}
        \lim_{x \to a} \frac{f'(x)}{g'(x)} = l
    \end{align*}
    allora esiste anche il limite
    \begin{align*}
        \lim_{x \to a} \frac{f(x)}{g(x)} = l
    \end{align*}
\end{theorem}
\end{document}
